% Generated. Do not edit.
% Projection scriptural-date-location/evidence, version 1: the Scriptural Date and
% Location page, generated whole from research/chronology.toml, the Scripture
% chronology corpus and the source library by tools/tpt document-types refresh.
% document: liturgy/roman-rite/1962/propers/temporal/50-tenth-after-pentecost
\section*{Appendix: Evidence for the Scriptural Dates}
\label{sec:date-location-evidence}
{\small For each passage, as the Scriptural Date and Location table names it: what is dated, how the corpus weighs each answer, the source's own words, and the sources cited. A subject the corpus does not date is reported as such and given no date.\par}
{\small
\par\smallskip\noindent\textit{Gradual: Ps. 16:8 (Vg.; Heb. 17), Ps. 16:2 (Vg.; Heb. 17)}\par\nobreak
\begin{itemize}[leftmargin=1.2em,itemsep=1pt,topsep=1pt,parsep=0pt]
\item \textbf{Historical setting}. No date: the chronology corpus has not yet dated this (research pending).
\item \textbf{Composition} (preferred). Psalms: Before c.~165 B.C. (``before the Maccabean period, around 165 B.C.; no individual psalm can be dated securely''; Catholic critical chronology; New American Bible, Revised Edition --- USCCB web state (2026), Psalms introduction).
\end{itemize}
\par\smallskip\noindent\textit{Offertory: Ps. 24:1--3 (Vg.; Heb. 25)}\par\nobreak
\begin{itemize}[leftmargin=1.2em,itemsep=1pt,topsep=1pt,parsep=0pt]
\item \textbf{Historical setting}. As under Gradual: Ps. 16:8 (Vg.; Heb. 17), Ps. 16:2 (Vg.; Heb. 17).
\item \textbf{Composition}. As under Gradual: Ps. 16:8 (Vg.; Heb. 17), Ps. 16:2 (Vg.; Heb. 17).
\end{itemize}
\par\smallskip\noindent\textit{Communion: Ps. 50:21 (Vg.; Heb. 51)}\par\nobreak
\begin{itemize}[leftmargin=1.2em,itemsep=1pt,topsep=1pt,parsep=0pt]
\item \textbf{Superscription setting} (preferred). David's sin with Bethsabee and his repentance under Nathan's rebuke: When Nathan the prophet came to him, after he had sinned with Bethsabee (The Holy Bible, Douay-Rheims (Challoner revision), Project Gutenberg eBook 1581: Ps. 50:1, 50:2, 2 Kings 12:13).
\item \textbf{Composition}. As under Gradual: Ps. 16:8 (Vg.; Heb. 17), Ps. 16:2 (Vg.; Heb. 17).
\end{itemize}
\par\smallskip\noindent\textit{Introit: Ps. 54:17--18, 20, 23 (Vg.; Heb. 55), Ps. 54:2 (Vg.; Heb. 55)}\par\nobreak
\begin{itemize}[leftmargin=1.2em,itemsep=1pt,topsep=1pt,parsep=0pt]
\item \textbf{Historical setting}. As under Gradual: Ps. 16:8 (Vg.; Heb. 17), Ps. 16:2 (Vg.; Heb. 17).
\item \textbf{Composition}. As under Gradual: Ps. 16:8 (Vg.; Heb. 17), Ps. 16:2 (Vg.; Heb. 17).
\end{itemize}
\par\smallskip\noindent\textit{Alleluia: Ps. 64:2 (Vg.; Heb. 65)}\par\nobreak
\begin{itemize}[leftmargin=1.2em,itemsep=1pt,topsep=1pt,parsep=0pt]
\item \textbf{Historical setting}. As under Gradual: Ps. 16:8 (Vg.; Heb. 17), Ps. 16:2 (Vg.; Heb. 17).
\item \textbf{Composition}. As under Gradual: Ps. 16:8 (Vg.; Heb. 17), Ps. 16:2 (Vg.; Heb. 17).
\end{itemize}
\par\smallskip\noindent\textit{Gospel: Luke 18:9--14}\par\nobreak
\begin{itemize}[leftmargin=1.2em,itemsep=1pt,topsep=1pt,parsep=0pt]
\item \textbf{Narrated event}. No date: the chronology corpus has not yet dated this (research pending).
\item \textbf{Composition} (disputed). The Gospel of St Luke: c.~A.D.~70 (``About the year 70''; The Catholic Encyclopedia, Volume 14 (New York, 1912)); Before the end of the Roman imprisonment, when the Acts was finished (De auctore, de tempore compositionis et de historica veritate Evangeliorum secundum Marcum et secundum Lucam (1912), responsa-i-ix).
\end{itemize}
\par\smallskip\noindent\textit{Epistle: 1 Cor. 12:2--11}\par\nobreak
\begin{itemize}[leftmargin=1.2em,itemsep=1pt,topsep=1pt,parsep=0pt]
\item \textbf{Composition} (disputed). The First Epistle to the Corinthians: A.D.~56 (``(1 and 2 Corinthians; Galatians), 56''; The Catholic Encyclopedia, Volume 11 (New York, 1911)); A.D.~57 (``about Easter A.D. 57''; The Catholic Encyclopedia, Volume 4 (New York, 1908)).
\end{itemize}
\par}
